\begin{table}[htbp]
\centering
\footnotesize
\caption{\textbf{2.2} Los d\'ias sin transacciones que \emph{no} son s\'abado. Los 53 s\'abados se omiten de la tabla por brevedad.}
\label{tab:p2_feriados}
\begin{tabular}{llll}
\toprule
\textbf{fecha} & \textbf{d\'ia} & \textbf{motivo} & \textbf{tipo} \\
\midrule
2010-12-24 & viernes & feriado bancario: Nochebuena (la tienda cierra & fijo \\
2010-12-26 & domingo & feriado bancario: Boxing Day & fijo \\
2010-12-27 & lunes & feriado bancario: Navidad (dia sustituto) & fijo \\
2010-12-28 & martes & feriado bancario: Boxing Day (dia sustituto) & fijo \\
2010-12-29 & miercoles & sin explicacion en el calendario & --- \\
2010-12-30 & jueves & sin explicacion en el calendario & --- \\
2010-12-31 & viernes & sin explicacion en el calendario & --- \\
2011-01-02 & domingo & sin explicacion en el calendario & --- \\
2011-01-03 & lunes & feriado bancario: Ano Nuevo (dia sustituto) & fijo \\
2011-04-22 & viernes & feriado bancario: Viernes Santo & \textbf{MOVIL} \\
2011-04-24 & domingo & sin explicacion en el calendario & --- \\
2011-04-25 & lunes & feriado bancario: Lunes de Pascua & \textbf{MOVIL} \\
2011-04-29 & viernes & feriado bancario: Boda real (feriado extraordi & \textbf{irrepetible} \\
2011-05-02 & lunes & feriado bancario: Early May Bank Holiday & fijo (primer lunes de mayo) \\
2011-05-30 & lunes & feriado bancario: Spring Bank Holiday & fijo (ultimo lunes de mayo) \\
2011-08-29 & lunes & feriado bancario: Summer Bank Holiday & fijo (ultimo lunes de agosto) \\
\bottomrule
\end{tabular}
\\[0.3em]{\scriptsize \textbf{Qu\'e leer:} tres de los feriados son m\'oviles o irrepetibles. La Pascua se mueve hasta un mes entre a\~nos y la boda real de 2011 no se repite nunca, as\'i que ning\'un modelo entrenado con un a\~no de datos puede anticiparlos.}
\end{table}
